% !TEX TS-program = lualatex

% FORCE-CROP

\documentclass[border = 3pt]{standalone}

\usepackage[svgnames]{xcolor}
\usepackage[3d]{luadraw}

\begin{document}

\begin{luadraw*}{name = isometric-prespective}
local ld = luadraw

local M      = ld.pt3d.M
local Origin = ld.pt3d.Origin
local vecI   = ld.pt3d.vecI
local vecJ   = ld.pt3d.vecJ
local vecK   = ld.pt3d.vecK

------
-- Ici, nous jouons plus finement avec le paramétrage de la
-- perspective centrale. Voici comment cela fonctionne.
--
--     + lua::''67'', lua::''25'' et lua::''42'' sont les ¨coords
--     cylindriques de l'observateur relativement au repère avec
--     pour nouvelle origine lua::''NewOrigin''.
--
--     + Le plan de projection passe par lua::''NewOrigin'' et
--     est orienté par le vecteur directeur allant du point
--     lua::''NewOrigin'' à l'observateur.
--
--
-- note::
--     En gardant le réglage lua::''window  = {-6, 6.5, -5, 5.5}'',
--     notre choix de positionnement, qui élève l'observateur,
--     fait disparaître la flèche de l'axe des ordonnées.
--     Choisir lua::''window  = {-6, 6.5, -10, 5.5}'' rectifie la
--     situation (mais cela devient moins pertinent pour expliquer
--     comment fonctionne le réglage de la perspective centrale).
--
--
-- warning::
--     La ¨dist cylindrique ne semble modifier que très peu le
--     rendu final.
--
--
-- note::
--     Choisir lua::''-90'' pour les deux angles donne une vue
--     à la ¨opengl avec l'axe des abscisses `x` vers la droite,
--     celui des côtes `z` vers le bas, et celui des ordonnées
--     `y` "sortant" du plan de la vue 2D.
------
local NewOrigin = Origin + 4*vecK

local graphview = ld.graph3d:new{
  viewdir = {"central", 67, 52, 42, NewOrigin},
  window  = {-6, 6.5, -5, 5.5},
  size    = {10, 10},
  bbox    = false,
}

graphview:Linewidth(4)
ld.Hiddenlinestyle = "dashed"

local P = ld.parallelep(
  Origin,
  4*vecI, 6*vecJ, 1.5*vecK
)

table.remove(P.facets,2)

P = ld.shift3d(
  P,
  - ld.pt3d.isobar3d(P.vertices)
)

graphview:Dscene3d(
  graphview:addPoly(
    P,
    {
      edge      = true,
      hidden    = true,
      contrast  = 0,
      twoside   = false,
      edgewidth = 6
    }
  ),
  graphview:addFacet(
    ld.getfacet(P,1) ,
    {color = "LightGray"}
  ),
  graphview:addAxes(
    Origin,
    {arrows = 1}
  )
)

graphview:Show()
\end{luadraw*}

\end{document}
